% PDF transcription; source and license in ../../tex/003_有限域椭圆曲线的Hasse界.tex
\paragraph{7.1 Separable and inseparable endomorphisms.}
Recall that the Frobenius endomorphism $\pi_E$ is inseparable. In order to prove Hasse's theorem we will need to use the fact that $\pi_E-1$ is separable. This follows from a much more general result: adding a separable isogeny to an inseparable isogeny always yields a separable isogeny. Note that the sum of two separable isogenies need not be separable: in characteristic $p>0$, if we have $a+b=p$ and both $a$ and $b$ prime to $p$, then $[a]$ and $[b]$ are both separable but $[a]+[b]=[a+b]=[p]$ is inseparable.

\begin{lemma}[7.1]
Let $\alpha$ and $\beta$ be isogenies from $E_1$ to $E_2$, with $\alpha$ inseparable. Then $\alpha+\beta$ is inseparable if and only if $\beta$ is inseparable.
\end{lemma}
\begin{proof}
If $\beta$ is inseparable then by Corollary 5.4 we can write $\alpha=\alpha_{\mathrm{sep}}\circ\pi^m$ and $\beta=\beta_{\mathrm{sep}}\circ\pi^n$, where $\pi$ is the $p$-power Frobenius map and $m,n>0$. We then have
\[
\alpha+\beta=\alpha_{\mathrm{sep}}\circ\pi^m+\beta_{\mathrm{sep}}\circ\pi^n
=(\alpha_{\mathrm{sep}}\circ\pi^{m-1}+\beta_{\mathrm{sep}}\circ\pi^{n-1})\circ\pi,
\]
which is inseparable (any composition involving an inseparable isogeny is inseparable because inseparable degrees multiply). If $\alpha+\beta$ is inseparable, then so is $-(\alpha+\beta)$, and $\alpha-(\alpha+\beta)=\beta$ is a sum of inseparable isogenies, which we have just shown is inseparable.
\end{proof}

\paragraph{Remark 7.2.}
Since the composition of an inseparable isogeny with any isogeny is always inseparable, Lemma 7.1 implies that the inseparable endomorphisms in $\End(E)$ form an ideal (provided we view $0$ as inseparable, which we do).

\paragraph{7.2 Hasse's Theorem.}
We are now ready to prove Hasse's theorem.
\begin{theorem}[7.3; Hasse]
Let $E/\mathbb F_q$ be an elliptic curve over a finite field. Then
\[
\#E(\mathbb F_q)=q+1-t,
\]
where $t:=\tr\pi_E$ is the trace of the Frobenius endomorphism $\pi_E$ and $|t|\le2\sqrt q$.
\end{theorem}
\begin{proof}
Recall that we defined $\mathbb F_q$ as the splitting field of $x^q-x$ over $\mathbb F_p$, where $p=\operatorname{char}(\mathbb F_q)$, thus
\[
\mathbb F_q=\{\alpha\in\overline{\mathbb F}_p:\alpha^q-\alpha=0\}
=\{\alpha\in\overline{\mathbb F}_q:\alpha^q-\alpha=0\}
\]
is precisely the subfield of $\overline{\mathbb F}_q$ fixed by the $q$-power Frobenius automorphism $x\mapsto x^q$. The Frobenius endomorphism $\pi_E:E\to E$ is defined by $\pi_E(x:y:z)=(x^q:y^q:z^q)$, therefore
\[
\begin{aligned}
E(\mathbb F_q)&=\{P\in E(\overline{\mathbb F}_q):\pi_E(P)=P\}\\
&=\{P\in E(\overline{\mathbb F}_q):\pi_E(P)-P=0\}=\ker(\pi_E-1),
\end{aligned}
\]
where $1$ denotes the multiplication-by-1 map $[1]\in\End(E)$. The Frobenius endomorphism $\pi_E$ is inseparable and $-1$ is separable, so by Lemma 7.1 the endomorphism $\pi_E-1$ is separable, thus the cardinality of its kernel is equal to its degree (by Theorem 5.8). Therefore
\[
\begin{aligned}
\#E(\mathbb F_q)&=\#\ker(\pi_E-1)=\deg(\pi_E-1)\\
&=(\widehat{\pi_E-1})(\pi_E-1)
=\widehat\pi_E\pi_E+1-(\widehat\pi_E+\pi_E)=q+1-t.
\end{aligned}
\]
It remains only to show that $|t|\le2\sqrt q$.

Consider the endomorphism $r\pi_E-s$ for $r,s\in\mathbb Z$ with $s\ne0$. We have
\[
\begin{aligned}
\deg(r\pi_E-s)&=(\widehat{r\pi_E-s})(r\pi_E-s)
=(\widehat\pi_E\widehat r-\widehat s)(r\pi_E-s)
=(\widehat\pi_Er-s)(r\pi_E-s)\\
&=\widehat\pi_Er^2\pi_E-\widehat\pi_Ers-sr\pi+s^2
=r^2\widehat\pi_E\pi_E-rs(\widehat\pi_E+\pi_E)+s^2\\
&=r^2\deg\pi_E-rs\tr\pi_E+s^2\\
&=r^2q-rst+s^2,
\end{aligned}
\]
where we have used Lemmas 6.11 and 6.12, and the fact that $\mathbb Z$ is in the center of $\End(E)$. Dividing by $s^2$ and noting that $\deg(r-\pi_Es)\ge0$ yields the inequality
\[
q(r/s)^2-t(r/s)+1\ge0,
\]
valid for all rational numbers $r/s$. Now $\mathbb Q$ is dense in $\mathbb R$, so we must have $qx^2-tx+1\ge0$ for all real numbers $x$. It follows that the discriminant $t^2-4q$ cannot be positive, which yields the desired bound $|t|\le2\sqrt q$.
\end{proof}
